\scanpage{268}
\mspage{80}
% archive sheet number: 268
\[
\operatorname{Hom}_{\unsure{k(f)|K(f)}}(x,y)\unsure{\text{---}|\longrightarrow}
\operatorname{Hom}_J(\bar f(x),\bar f(y))
\]

Si $x,y\in\mathrm{Ob}\,J$, c'est l'identité, si $x\unsure{,}y\in I$, c'est $\bar f$. Si $x\in I$, $y\in J$, on doit définir une flèche

\[
\operatorname{Hom}(f(x),y)\longrightarrow\operatorname{Hom}(f(x),y)
\]

et on prend encore l'identité. \deleted{Le foncteur} On a

\diagram{\begin{tikzcd}[ampersand replacement=\&,column sep=small,row sep=large]
J\arrow[r,hook,"j"]\arrow[dr,equal]\&K(f)\arrow[d,"\bar f"]\& I\arrow[l,hook,no head,"i"']\\
\&J\&
\end{tikzcd}}

\[
\left\{\begin{aligned}
\bar f\circ j&=\mathrm{id}_J\\
j\circ\bar f&\overset{\deleted{\unsure{?}}}{=}
\begin{cases}\deleted{\unsure{*}}j&\text{sur }J\\jf&\text{sur }I\end{cases}
\end{aligned}\right.
\qquad\text{i.e.}\quad
\underset{\deleted{\text{sur }I}}{\boxed{\begin{aligned}
j\bar f j&=j\ \text{sur }J\\
j\bar f i&=jf\ \text{sur }I
\end{aligned}}}
\]

et on a \deleted{\illegible{1}}

\[
\begin{aligned}
\alpha:\mathrm{id}_K&\longrightarrow j\bar f&&\text{i.e. hom fonctoriel}\\
x&\overset{\alpha(\unsure{x|0})}{\longmapsto}j(\bar f(x))&&x\in\mathrm{Ob}\,K
\end{aligned}
\]

défini ainsi :

\[
\left\{\begin{aligned}
&\text{si }x\in J,\quad\alpha(x)=\mathrm{id}_x,\ \text{\unsure{s.à|s}}:\quad
 x\xrightarrow{\sim}j(\bar f(x))=j(x)=\deleted{\unsure{j}}x\\
&\text{si }x\in I\qquad\alpha(x)\mathrel{\unsure{\doteq|=}} x\underset{\text{ds }K(x)}{\longrightarrow}f(x),\ \text{ou }
 \underbrace{f(x)\longrightarrow f(x)}_{\text{ds }I}\\
&\hspace{18em}\alpha(x)=\mathrm{id}_{f(x)},
\end{aligned}\right.
\]

On doit vérifier la \deleted{\unsure{compatibilité}}\inserted{fonctorialité} pour $x$ variable dans $K$\ldots

Soit $\Phi$ un $\mathcal M$-préfaisceau sur $K$. Il définit

\[
\begin{aligned}
F&\mathrel{\unsure{\doteq|=}}i^*(\Phi)&&\in\mathrm{Ob}\,\mathbb D_{\mathcal M}(I)\\
G&=j^*(\Phi)&&\in\mathrm{Ob}\,\mathbb D_{\mathcal M}(J)
\end{aligned}
\]

D'autre part $f^*(G)$ s'interprète comme $(j\bar f)^*(\Phi)|I$ car

\[
\begin{aligned}
f^*(G)(i)&=G(\underset{\bar f(i)\,\unsure{=|\parallel}\ \in J}{\underline{f(i)}})
 =j^*(\Phi)(\bar f(i))=\deleted{\unsure{G(f(i))}}\Phi(j\bar f(i))\\
&=(j\bar f)^*(\Phi)(i)
\end{aligned}
\]
